\documentclass[12pt, letterpaper]{article}


\usepackage[utf8]{inputenc}
\usepackage[spanish]{babel}


\usepackage[margin=2.54cm]{geometry}


\usepackage{times}
\usepackage{setspace}
\doublespacing 

\usepackage{hyperref}
\hypersetup{
    colorlinks=true,
    linkcolor=blue,
    filecolor=magenta,      
    urlcolor=blue,
}

\begin{document}

\begin{titlepage}
    \centering
    \vspace*{2cm}
    
    {\Large \bfseries El Buen Uso de la Inteligencia Artificial en el Ámbito Estudiantil\par}
    
    \vspace{2cm}
    
    {\large DANIEL ALEXANDER MIRANDA CASTILLO \par}
    \vspace{0.5cm}
    {\large Universidad Mariano Galvez \par}
    {\large SEMINARIO DE TECNOLOGÍAS DE INFORMACIÓN \par}
    {\large Mélvin Cali\par}
    \vspace{1cm}
    {\large 24 de julio de 2026 \par}
    
    \vfill
\end{titlepage}

\section*{Introducción}
La inteligencia artificial está avanzando a pasos agigantados y, claro, ya se metió de lleno en la educación. La verdad es una herramienta que sirve un montón para buscar información o para que los estudiantes aprendan por su cuenta, pero también hay que verle el otro lado. Su uso sin control está trayendo un pisto de problemas éticos y pedagógicos que no podemos tapar con la mano. Como ingenieros y docentes, nos toca analizar bien cómo los jovenes están usando estas tecnologías, pero para hacerlo bien necesitamos metodologías y herramientas de recolección de datos que sean claras, objetivas y que realmente reflejen lo que pasa en las aulas de nuestro país, no solo teorías importadas.

\section*{Criterios para el Diseño de una Encuesta Efectiva}
Para llevar a cabo una investigación cuantitativa sólida sobre el buen uso de la IA en el ámbito estudiantil, la encuesta utilizada debe cumplir con los siguientes principios metodológicos:

\begin{itemize}
    \item \textbf{Objetivo delimitado:} Definir con precisión si se investiga la ética, la frecuencia de uso o el impacto del uso de IA en el rendimiento académico.
    \item \textbf{Brevedad y precisión:} Redactar preguntas claras que no excedan un tiempo de respuesta de 3 a 5 minutos, garantizando una mayor participación.
    \item \textbf{Neutralidad en la redacción:} Formulaciones impersonales para evitar sesgar o juzgar la honestidad académica del encuestado.
   
\end{itemize}

\section*{Aplicación Práctica: Enlace a la Encuesta}
Con base en los criterios metodológicos investigados, se diseñó una encuesta breve para diagnosticar las prácticas éticas y el uso responsable de la IA por parte de los estudiantes. Se puede acceder a dicho instrumento en línea a través del siguiente enlace:

\begin{center}
    \url{https://docs.google.com/forms/d/e/1FAIpQLSfWryjtD0TNV8cxtZohIiWQwwVl4Yebo1CGno2XJcetPXdTtA/viewform?usp=publish-editor}
\end{center}

\section*{Conclusión}
El buen uso de la Inteligencia Artificial por parte de los estudiantes no se trata de prohibirla, sino de usarla como una herramienta que potencie el pensamiento crítico, no que lo reemplace. Porque, seamos sinceros, la IA es un apoyo, pero el esfuerzo mental y el análisis siguen siendo responsabilidad del estudiante. Para medir cómo están usando estas tecnologías en las aulas guatemaltecas, lo ideal es diseñar encuestas bien hechas, claras y con preguntas que de verdad sirvan. Solo así las universidades y escuelas pueden tomar decisiones acertadas sobre cómo regular su uso y cómo enseñar a los estudiantes a usarlas con cabeza

\newpage

\section*{Referencias y Código Fuente}

\subsection*{Referencias}
\begin{itemize}
    \item American Psychological Association. (2020). \textit{Publication manual of the American Psychological Association} (7th ed.). https://doi.org/10.1037/0000165-000
    \item Hernández-Sampieri, R., \& Mendoza, C. P. (2018). \textit{Metodología de la investigación: Las rutas cuantitativa, cualitativa y mixta}. McGraw-Hill Education.
\end{itemize}

\subsection*{Repositorio del Código Fuente (Git)}
El código fuente en LaTeX utilizado para la creación y formato de este documento ubicado en:

\begin{center}
    \url{https://github.com/dmirandac5-tech/Foro-acad-mico-1.git}
\end{center}

\end{document}